\section{Laboratorios del curso y su aporte al proyecto}

\begin{table}[H]
\centering
\small
\caption{Recorrido de los laboratorios y su producto}
\begin{tabular}{@{}p{3.6cm}p{5.4cm}p{4.6cm}@{}}
\toprule
\textbf{Laboratorio} & \textbf{Qué se hizo} & \textbf{Aporte al proyecto} \\
\midrule
Python científico reproducible & Simulación y análisis de señales: estadística,
cálculo numérico, álgebra lineal, visualización; informe y presentación en LaTeX &
Entorno con \texttt{uv}, Ruff y pruebas; paquete importable; reproducción en otro
computador \\
\addlinespace
Adquisición y cuantización & Captura controlada de imágenes, cuantización uniforme
y análisis de espacios de color & Criterios de cuantización y de espacios de color
usados en la caracterización \\
\addlinespace
Formulación del proyecto & Problema, pregunta, objetivos, selección de datos,
matriz de técnicas, cronograma y riesgos & Sección 1 y la matriz del anexo B \\
\addlinespace
Caracterización (Fase 1) & Manifiesto, verificación de integridad, calidad y
partición por paciente & Sección 2 \\
\addlinespace
Análisis exploratorio y tabla de características & Campo irradiado, máscaras de
referencia, nueve estrategias, descriptores y tabla & Secciones 3 y 4 (OE3, OE4) \\
\addlinespace
Modelado supervisado & Regresión múltiple con protocolo congelado y práctica común
con regresión logística & Sección 4, modelado \\
\addlinespace
Fase 3 & OE2 en cinco lotes, OE3 en validación, OE4, OE5 y bootstrap & Sección 4
completa \\
\bottomrule
\end{tabular}
\end{table}